\documentclass[11pt]{article}
\usepackage{amsmath,amssymb,amsthm,geometry,longtable,hyperref}
\geometry{margin=24mm}
\newcommand{\D}[1]{|D_{19}^{(#1)}\rangle}
\newcommand{\Tr}{\operatorname{Tr}}
\newtheorem{theorem}{Theorem}
\title{An exact algebraic four-dimensional second-order anticoherent subspace at nineteen qubits}
\author{Catalogue calculation}\date{30 October 2026}
\begin{document}\maketitle
The construction below is an exact isolated algebraic root definition with an explicit
rational existence certificate. It does not assert a short formula in radicals.
Every vector in the resulting four-dimensional complex span is second-order AC.

\section{Compact analytical basis}
For $j=0,\ldots,4$, define
\[
 a_j=\sqrt{\binom{19}{4j}}(A_j+iB_j),\qquad
 b_j=\sqrt{\binom{19}{4j+1}}(C_j+iD_j),\qquad B_0=D_0=0,
\]
where the eighteen real coordinates are ordered as
\begin{align*}
(A_0,A_1,B_1,A_2,B_2,A_3,B_3,A_4,B_4)&=(x_1,\ldots,x_9),\\
(C_0,C_1,D_1,C_2,D_2,C_3,D_3,C_4,D_4)&=(x_{10},\ldots,x_{18}).
\end{align*}
The exact definition of $x$ is given in Sections 2--3. Set
\begin{align*}
|u_1\rangle&=\sum_{j=0}^4a_j\D{4j},&
|u_2\rangle&=\sum_{j=0}^4b_j\D{4j+1},\\
|u_3\rangle&=\sum_{j=0}^4b_j\D{18-4j},&
|u_4\rangle&=\sum_{j=0}^4a_j\D{19-4j}.
\end{align*}
These four vectors have disjoint residue-class supports modulo four.
The reversal is unitary and contains no coefficientwise complex conjugation.

\begin{theorem}
The specified rational box contains a unique zero of the eighteen quadratic
polynomials defined below. Its coordinates yield an orthonormal basis satisfying
\[
R_2(|u_a\rangle\langle u_b|)=\delta_{ab}\frac{I_3}{3},\qquad 1\le a,b\le4.
\]
Consequently the span is second-order AC, and $K(19,2)=4$.
\end{theorem}

\section{The exact quadratic equations}
Write $c_{a,r}=\sqrt{\binom{19}{r}}\,\xi_{a,r}$ for the displayed coefficients.
Thus $\xi_{1,4j}=A_j+iB_j$, $\xi_{2,4j+1}=C_j+iD_j$,
$\xi_{3,18-4j}=C_j+iD_j$ and $\xi_{4,19-4j}=A_j+iB_j$;
all other entries vanish. Define
\[
h_0(r)=1,\quad h_1(r)=19-2r,\quad
h_2(r)=3(19-2r)^2-399.
\]
For the triples $(q,p)$ used in the table, put
\[
P_{ab}^{q,p}(x)=
\sum_{r=q}^{19}\binom{19}{r}\,r^{\underline q}h_p(r)
\overline{\xi_{a,r-q}}\xi_{b,r}
-\delta_{q0}\delta_{p0}\delta_{ab}.
\]
Here $r^{\underline q}=r(r-1)\cdots(r-q+1)$, with value one for $q=0$.
The equation $F_i=0$ is the listed real or imaginary part of $P_{ab}^{q,p}$:
\[
\begin{array}{c|ccccc}
i&a&b&q&p&\text{part}\\\hline
1 & 1 & 1 & 0 & 0 & \mathrm{real} \\
2 & 2 & 2 & 0 & 0 & \mathrm{real} \\
3 & 1 & 1 & 0 & 1 & \mathrm{real} \\
4 & 2 & 2 & 0 & 1 & \mathrm{real} \\
5 & 1 & 1 & 0 & 2 & \mathrm{real} \\
6 & 2 & 2 & 0 & 2 & \mathrm{real} \\
7 & 1 & 2 & 1 & 0 & \mathrm{real} \\
8 & 1 & 2 & 1 & 0 & \mathrm{imaginary} \\
9 & 2 & 3 & 1 & 0 & \mathrm{real} \\
10 & 4 & 1 & 1 & 0 & \mathrm{real} \\
11 & 1 & 2 & 1 & 1 & \mathrm{real} \\
12 & 1 & 2 & 1 & 1 & \mathrm{imaginary} \\
13 & 2 & 3 & 1 & 1 & \mathrm{imaginary} \\
14 & 4 & 1 & 1 & 1 & \mathrm{imaginary} \\
15 & 1 & 3 & 2 & 0 & \mathrm{real} \\
16 & 1 & 3 & 2 & 0 & \mathrm{imaginary} \\
17 & 3 & 1 & 2 & 0 & \mathrm{real} \\
18 & 3 & 1 & 2 & 0 & \mathrm{imaginary} \\
\end{array}
\]
These are explicitly specified quadratic polynomials with integer coefficients.
The factor $\binom{19}{r}r^{\underline q}$ follows directly from normalized Dicke
matrix elements of $J_+^q$.

For $q=0$ and $p=0,1,2$, the resulting equations are respectively the Gram,
$2J_z$ and $12J_z^2-399I$ compression conditions. The remaining cases are
$(q,p)=(1,0),(1,1),(2,0)$, giving $J_+$, $2J_+J_z$ and $J_+^2$.
Their adjoints give the lowering conditions. Together these are exactly
normalization and vanishing tensor ranks one and two.

Before linear redundancy removal, the prescribed support and reversal leave
34 nonzero real equations. The certificate contains a rational $34\times18$
reconstruction matrix expressing each original polynomial as a linear
combination of the selected eighteen. Both the physical polynomial coefficients
and this reconstruction are checked exactly. The independent verifier instead
reconstructs every reduced-state polynomial directly from Dicke splitting and
checks all 144 complex entries, including every off-diagonal basis pair.

\section{Rational box and existence proof}
Set $r_0=10^{-50}$ and $x_i^{(0)}=n_i/10^{70}$, where the integers $n_i$
are listed in the appendix. The exact coordinate definition is the unique root
of $F(x)=0$ in
\[
\mathcal B=\{x\in\mathbb R^{18}:|x_i-x_i^{(0)}|\le r_0\text{ for all }i\}.
\]
The certificate supplies an explicit rational matrix $A$ with denominator
$10^{60}$. Integer arithmetic verifies, in the induced infinity norm,
\begin{align*}
\delta=\|I-AF'(x^{(0)})\|_\infty&<2\cdot10^{-55},\\
\eta=\|AF(x^{(0)})\|_\infty&<5\cdot10^{-71},\\
\|A\|_\infty&<3/10.
\end{align*}
For every $x\in\mathcal B$, the exact polynomial Hessians give
\[
\|F'(x)-F'(x^{(0)})\|_\infty
\le181138248\,\|x-x^{(0)}\|_\infty.
\]
The map $G(x)=x-AF(x)$ is therefore a strict contraction on $\mathcal B$, with
\[
\operatorname{Lip}(G)<7\cdot10^{-43},\qquad
\|G(x)-x^{(0)}\|_\infty<6\cdot10^{-71}<r_0.
\]
It maps the box strictly into itself. The contraction theorem supplies a unique
fixed point. Since $\delta<1$, the square matrix $AF'(x^{(0)})$ and hence $A$
are invertible. The fixed point consequently satisfies $F(x)=0$.
Every coordinate is real algebraic: the invertible Jacobian makes this an
isolated nonsingular root of an integer polynomial system. Decimal optimization
was used only to locate the rational box; none of the proof inequalities relies
on floating-point arithmetic.

\section{Whole-span identity and optimal dimension}
The moment conditions are equivalent to the full two-qubit cross identity.
For a direct expression, define
\[
M(c)_{ab}=c_{a+b}\sqrt{\frac{\binom2a\binom{17}b}{\binom{19}{a+b}}},
\quad0\le a\le2,\quad0\le b\le17.
\]
The independent certificate verifies $M(u_i)M(u_j)^\dagger=\delta_{ij}I_3/3$.
Linearity now proves the identity for every normalized complex superposition.

For the upper bound, couple physical spin $j=19/2$ to auxiliary spin $3/2$ and
write $X=\mathbf J\cdot\mathbf S$. The second-order AC conditions give zero
compression of
\[
W=X^2-\frac{15}{2}X-\frac{1995}{16}I
\]
on the $4k$-dimensional code--auxiliary space. On the four total-spin sectors
$8,9,10,11$, only the spin-eight sector is positive; it has dimension 17, and
there are no zero eigenvalues. Projection of an isotropic subspace into the
positive eigenspace is injective. Thus $4k\le17$, giving $k\le4$.
The displayed basis attains this bound, proving $K(19,2)=4$ without restricting
the symmetry of a competing subspace.

\section{Relation to the numerical source and Majorana points}
The source is \texttt{\(\Pi\)\_19\_4\_2} in the supplied \texttt{ACLiege.jl},
line 1006. Its span has approximately one vector in each residue class modulo
four. A common coefficient rephasing $c_r\mapsto e^{i\beta r}c_r$,
$\beta\simeq0.020395581511200894$, makes the paired residue vectors related by
unitary reversal up to individual phases. Refinement of this structure produces
the isolated algebraic root above. The refined projector differs from the
rotated source projector by $3.76\cdot10^{-8}$ in Frobenius norm. This is an
approximate association with the supplied decimal subspace, not an assertion
that the rounded source already satisfies exact identities.

The four Majorana polynomials reduce to quartics in $z^4$ times an explicit
power of $z$. Exact rational interval estimates and disjoint Rouch\'e disks
certify all sixteen reduced roots as simple and nonzero. The multiplicity
patterns are $(3,1^{16})$ for $u_1,u_4$ and $(2,1^{17})$ for $u_2,u_3$.
Their numerical coordinates reconstruct the state rays to less than
$1.2\cdot10^{-15}$. These are properties of the chosen basis representatives,
not constellations of the full four-dimensional subspace.

\appendix\section{Exact center numerators}
All entries have common denominator $10^{70}$. They define the box, not the
algebraic root itself.
{\scriptsize
\begin{longtable}{r|r}
i&n_i\\\hline
1 & 4291256405878055156282314634547861427979220094757699665120090204232617 \\
2 & -33403743290029659862230629603920758886371636037308289756601292022702 \\
3 & 9286452461006488056330409251798418783587222115934851191894963828098 \\
4 & -5225798200773107665673406364103769277465079282511927315613140169612 \\
5 & -19301092704014119073796814747460504928080621569302022539128686709246 \\
6 & -16815208565591816582138855416319178081433842169200896779507296881748 \\
7 & -1859347921660105387171264520156560634372276895955762161928623414220 \\
8 & 102508475707419101112407734202129538007790953632166242317018315249144 \\
9 & -151024968462427430345361080359081266605406391146222878545114079643221 \\
10 & 971387704479272983962329833995547339634441471688729989604032800569841 \\
11 & -28057426574450962205245346076324857905998642734238630969102319918354 \\
12 & -33824453079610240780609723472867009655839610826623582844498397024739 \\
13 & -5962445703294221169600092178690019470723325501597045975184766308450 \\
14 & 9136760453529342792627321184905574395292792392759686855533982183605 \\
15 & 31179359891046315480697005725986116066482343969177461969109788812489 \\
16 & -1312989377740439525583823548920649516725029788244065264526473257950 \\
17 & 241420505162928586308905175120113510060537141983884608679991329098025 \\
18 & -267419905258932840961826849307769699123163320347564966574538171275315 \\
\end{longtable}}
The machine-readable files \texttt{basis.json} and \texttt{certificate.json}
contain the same expressions, intervals, complete polynomials, reconstruction
matrix and rational approximate inverse. The two verifiers use independently
constructed moment and reduced-state polynomial systems.
\end{document}
